%  DESARROLLO PASO A PASO — EJERCICIO 2.3
%  Inecuación con raíz anidada: sqrt(x - sqrt(2-x)) < 1
%  Usach Premium — Cálculo I
%  Compila con: pdflatex c1_desarrollo_23.tex  (2 veces)
% ============================================================

\documentclass[letterpaper, 12pt]{article}

% ============================================================
% PAQUETES
% ============================================================
\usepackage[utf8]{inputenc}
\usepackage[spanish]{babel}
\usepackage{amsmath, amssymb}
\usepackage{geometry}
\usepackage{fancyhdr}
\usepackage{enumitem}
\usepackage{graphicx}
\usepackage{hyperref}
\usepackage{parskip}
\usepackage{xcolor}
\usepackage{tcolorbox}
\usepackage{eso-pic}
\usepackage{tikz}
\usepackage{array}
\usepackage{booktabs}
\usepackage{colortbl}
\usetikzlibrary{patterns, arrows.meta, decorations.pathmorphing, babel}
\tcbuselibrary{skins, breakable}

% ============================================================
% GEOMETRÍA
% ============================================================
\geometry{letterpaper, margin=2cm, top=3cm, bottom=2.5cm, headheight=2cm}

% ============================================================
% COLORES INSTITUCIONALES
% ============================================================
\definecolor{celesteSuave}{HTML}{F0F7FF}
\definecolor{azulFuerte}{HTML}{1A5276}
\definecolor{verdePaso}{HTML}{1E8449}
\definecolor{rojoAlerta}{HTML}{C0392B}
\definecolor{naranjaAcento}{HTML}{D35400}
\definecolor{grisTabla}{HTML}{D5D8DC}
\definecolor{enlace}{HTML}{1A5276}

% ============================================================
% RUTAS DE IMÁGENES
% ============================================================
\newcommand{\logoup}{../../../../../template/images/logo_up.png}
\newcommand{\logousach}{../../../../../template/images/logo_usach.png}
\newcommand{\logofondo}{../../../../../template/images/logo_up.png}

% ============================================================
% INFORMACIÓN
% ============================================================
\newcommand{\institucion}{Usach Premium}
\newcommand{\curso}{Cálculo I}
\newcommand{\guiasnumero}{Desarrollo — Ejercicio 2.3}
\newcommand{\semestre}{1er. Semestre de 2026}
\newcommand{\preparadopor}{Usach Premium.}
\newcommand{\webinstitucion}{www.usachpremium.cl}
\newcommand{\instagraminstitucion}{https://www.instagram.com/usach.premium}
\newcommand{\transparencia}{0.04}
\newcommand{\fraseportada}{"Por y para estudiantes"}

% ============================================================
% HYPERREF
% ============================================================
\hypersetup{colorlinks=true, linkcolor=enlace, urlcolor=enlace,
            citecolor=enlace, pdfborder={0 0 0}}

% ============================================================
% ENCABEZADO Y PIE
% ============================================================
\pagestyle{fancy}
\fancyhf{}
\fancyhead[L]{\includegraphics[height=1.2cm]{\logoup}}
\fancyhead[R]{%
    \begin{minipage}[b]{0.45\textwidth}
        \raggedleft\fontsize{9pt}{11pt}\selectfont
        \textbf{\color{azulFuerte}\institucion} \\
        \curso\ — \guiasnumero \\
        \textit{\color{gray}@usach.premium}
    \end{minipage}\hspace{0.1cm}%
    \includegraphics[height=1.2cm]{\logousach}}
\fancyfoot[L]{\fontsize{9pt}{11pt}\selectfont \textit{\fraseportada}}
\fancyfoot[R]{\fontsize{9pt}{11pt}\selectfont Página \thepage}
\renewcommand{\headrulewidth}{0.4pt}
\renewcommand{\footrulewidth}{0.4pt}
\fancypagestyle{plain}{\fancyhf{}\renewcommand{\headrulewidth}{0pt}
                                   \renewcommand{\footrulewidth}{0pt}}

% ============================================================
% MARCA DE AGUA
% ============================================================
\newcommand{\MarcaAgua}{
    \AddToShipoutPicture{
        \AtPageCenter{
            \begin{tikzpicture}[remember picture, overlay]
                \node[opacity=\transparencia] at (0,0)
                    {\includegraphics[width=0.7\textwidth]{\logofondo}};
            \end{tikzpicture}}}}

% ============================================================
% CAJAS
% ============================================================
% Caja de paso numerado
\newtcolorbox{paso}[2][]{
    colback=celesteSuave, colframe=azulFuerte, fonttitle=\bfseries,
    coltitle=white, colbacktitle=azulFuerte,
    title={\large Paso #2},
    arc=6pt, boxrule=1pt, enhanced, #1
}

% Caja de resultado/conclusión
\newtcolorbox{resultado}{
    colback=azulFuerte!10, colframe=azulFuerte, fonttitle=\bfseries,
    coltitle=white, colbacktitle=azulFuerte,
    title={\large\faCheck\ Resultado Final},
    arc=8pt, boxrule=1.5pt, enhanced,
    drop shadow={black!10}
}

% Caja de verificación Python
\newtcolorbox{verificacion}{
    colback=verdePaso!8, colframe=verdePaso, fonttitle=\bfseries,
    coltitle=white, colbacktitle=verdePaso,
    title={Verificación con Python (SymPy + NumPy)},
    arc=6pt, boxrule=1pt, enhanced,
    fontupper=\small\ttfamily
}

% Caja de advertencia/justificación
\newtcolorbox{justif}[1]{
    colback=naranjaAcento!8, colframe=naranjaAcento, fonttitle=\bfseries\small,
    title={#1}, arc=5pt, boxrule=0.8pt, enhanced
}

% ============================================================
\begin{document}
\thispagestyle{plain}
\MarcaAgua

% ============================================================
% PORTADA
% ============================================================
\AddToShipoutPicture*{%
    \put(54,700){\includegraphics[width=2cm]{\logoup}}
    \put(420,706){\includegraphics[width=5cm]{\logousach}}}

\vspace*{0.5cm}
\begin{center}
    \color{azulFuerte}
    {\huge \textbf{DESARROLLO PASO A PASO}} \\[0.5cm]
    {\Huge \textbf{\curso}} \\[0.4cm]
    {\Large Ejercicio 2.3 — Inecuación con raíz anidada} \\[1.2cm]
    \begin{tcolorbox}[colback=celesteSuave, colframe=azulFuerte, arc=10pt,
        width=0.88\textwidth, center, boxrule=1.5pt]
        \centering\vspace{0.3cm}
        {\large \textbf{Inecuación a resolver}}\\[0.6cm]
        \color{black}
        {\LARGE $\displaystyle \sqrt{x - \sqrt{2-x}} \;<\; 1$}
        \vspace{0.4cm}
    \end{tcolorbox}
\end{center}

\vspace{1.2cm}
\begin{center}
    {\Large \textbf{\semestre}} \\[0.7cm]
    {\large \textbf{\preparadopor}}
\end{center}

\vspace{1.8cm}
\begin{center}
    {\color{azulFuerte}\hrule height 0.5pt}\vspace{0.3cm}
    \begin{minipage}{0.45\textwidth}\centering
        \textbf{Instagram:} \\\small\href{\instagraminstitucion}{@usach.premium}
    \end{minipage}
    \begin{minipage}{0.45\textwidth}\centering
        \textbf{Web:} \\\small\href{http://\webinstitucion}{\webinstitucion}
    \end{minipage}
\end{center}

\pagenumbering{arabic}
\newpage

% ============================================================
% ENUNCIADO
% ============================================================
\begin{tcolorbox}[colback=azulFuerte!12, colframe=azulFuerte, arc=8pt,
    boxrule=1.2pt, enhanced]
    \textbf{\large Ejercicio 2.3.} \quad Resuelva la siguiente inecuación:
    \[
        \sqrt{x - \sqrt{2-x}} \;<\; 1
    \]
\end{tcolorbox}

\bigskip

% ============================================================
% PASO 1 — DOMINIO
% ============================================================
\begin{paso}{1: Determinación del Dominio}

Para que la expresión esté definida en $\mathbb{R}$ deben cumplirse
\textbf{simultáneamente} las siguientes condiciones:

\medskip
\textbf{(i) Raíz interior:} $\;2 - x \geq 0$

\[
    2 - x \geq 0 \quad\Longrightarrow\quad x \leq 2
\]

\textbf{(ii) Raíz exterior:} $\;x - \sqrt{2-x} \geq 0$

\[
    x \geq \sqrt{2-x}
\]

Como $\sqrt{2-x} \geq 0$, esta condición sólo puede satisfacerse para $x \geq 0$.
Elevando al cuadrado ambos lados (válido pues ambos son no negativos):

\[
    x^{2} \geq 2 - x \quad\Longrightarrow\quad
    x^{2} + x - 2 \geq 0 \quad\Longrightarrow\quad
    (x+2)(x-1) \geq 0
\]

\begin{center}
\renewcommand{\arraystretch}{1.4}
\begin{tabular}{c|ccccc}
    & $x < -2$ & $x=-2$ & $-2 < x < 1$ & $x=1$ & $x > 1$ \\
\hline
$(x+2)$ & $-$ & $0$ & $+$ & $+$ & $+$ \\
$(x-1)$ & $-$ & $-$ & $-$ & $0$ & $+$ \\
\hline
$(x+2)(x-1)$ & $+$ & $0$ & $-$ & $0$ & $+$ \\
\end{tabular}
\end{center}

\medskip
$(x+2)(x-1) \geq 0 \;\Rightarrow\; x \leq -2 \;\cup\; x \geq 1$.
Combinando con $x \geq 0$: $\quad x \geq 1$.

\bigskip
\textbf{Intersección de condiciones (i) y (ii):}
\[
    \{x \leq 2\} \;\cap\; \{x \geq 1\} \;=\; \boxed{Dom\,f = [\,1,\,2\,]}
\]

\end{paso}

\bigskip

% ============================================================
% PASO 2 — PRIMERA ELEVACIÓN AL CUADRADO
% ============================================================
\begin{paso}{2: Primera elevación al cuadrado}

Sobre el dominio $[1,2]$, el lado izquierdo $\sqrt{x-\sqrt{2-x}} \geq 0$
y el lado derecho $1 > 0$, luego la inecuación es equivalente a:

\begin{justif}{¿Por qué podemos elevar al cuadrado?}
    Si $0 \leq A < B$, entonces $A^2 < B^2$. Aquí $A = \sqrt{x-\sqrt{2-x}}
    \geq 0$ y $B = 1 > 0$, por lo que la equivalencia es válida y
    \textbf{no cambia el sentido de la desigualdad}.
\end{justif}

\[
    \sqrt{x - \sqrt{2-x}} < 1 \quad\stackrel{(\,)^2}{\Longleftrightarrow}\quad
    x - \sqrt{2-x} < 1
\]

Reordenando:
\[
    \boxed{x - 1 < \sqrt{2-x}}
\]

\end{paso}

\bigskip

% ============================================================
% PASO 3 — SEGUNDA ELEVACIÓN AL CUADRADO
% ============================================================
\begin{paso}{3: Segunda elevación al cuadrado}

En el dominio $[1,2]$ se tiene $x - 1 \geq 0$ y $\sqrt{2-x} \geq 0$,
por lo que podemos elevar al cuadrado nuevamente sin cambiar la desigualdad:

\begin{justif}{Justificación}
    $x-1 \geq 0$ en $[1,2]$ y $\sqrt{2-x}\geq 0$ siempre, luego
    $A = x-1 \geq 0$, $B = \sqrt{2-x} \geq 0$, y $A < B
    \Leftrightarrow A^2 < B^2$.
\end{justif}

\[
    x - 1 < \sqrt{2-x} \quad\stackrel{(\,)^2}{\Longleftrightarrow}\quad
    (x-1)^2 < 2-x
\]

Expandiendo y simplificando:
\[
    x^2 - 2x + 1 < 2 - x
\]
\[
    x^2 - 2x + 1 - 2 + x < 0
\]
\[
    \boxed{x^2 - x - 1 < 0}
\]

\textbf{Raíces del trinomio} $x^2 - x - 1 = 0$ usando la fórmula cuadrática:
\[
    x = \frac{1 \pm \sqrt{1 + 4}}{2} = \frac{1 \pm \sqrt{5}}{2}
\]
\[
    r_1 = \frac{1 - \sqrt{5}}{2} \approx -0.618 \qquad
    r_2 = \frac{1 + \sqrt{5}}{2} \approx 1.618
\]

\end{paso}

\newpage

% ============================================================
% PASO 4 — TABLA DE SIGNOS
% ============================================================
\begin{paso}{4: Tabla de signos de $q(x) = x^2 - x - 1$}

Factorizando: $q(x) = \left(x - \dfrac{1-\sqrt{5}}{2}\right)\left(x - \dfrac{1+\sqrt{5}}{2}\right)$

\medskip
\begin{center}
\renewcommand{\arraystretch}{1.7}
\begin{tabular}{c|ccccc}
    & \small$x < r_1$ & \small$x = r_1$ & \small$r_1 < x < r_2$ & \small$x = r_2$ & \small$x > r_2$ \\
\hline
$\left(x - r_1\right)$ & $-$ & $0$ & $+$ & $+$ & $+$ \\
$\left(x - r_2\right)$ & $-$ & $-$ & $-$ & $0$ & $+$ \\
\hline
\rowcolor{celesteSuave}
$q(x) = x^2 - x - 1$ & $+$ & $0$ & $\mathbf{-}$ & $0$ & $+$ \\
\end{tabular}
\end{center}

\medskip
donde $r_1 = \dfrac{1-\sqrt{5}}{2}$ y $r_2 = \dfrac{1+\sqrt{5}}{2}$.

\bigskip
\textbf{Recta numérica:}

\begin{center}
\begin{tikzpicture}[scale=1.1]
    \draw[->] (-2.5,0) -- (4,0) node[right] {$x$};
    % raices
    \fill[azulFuerte] (-0.618,0) circle (2.5pt) node[below=4pt] {\small$r_1{\approx}-0.618$};
    \fill[azulFuerte]  (1.618,0) circle (2.5pt) node[below=4pt] {\small$r_2{\approx}1.618$};
    % signos
    \node at (-1.6, 0.35) {\color{rojoAlerta}$q > 0$};
    \node at ( 0.5, 0.35) {\color{verdePaso}$q < 0$};
    \node at ( 2.8, 0.35) {\color{rojoAlerta}$q > 0$};
    % arco negativo
    \draw[verdePaso, very thick] (-0.618,0) -- (1.618,0);
    % dominio [1,2]
    \draw[azulFuerte!60, dashed, thick] (1,-0.55) -- (1,0.55);
    \draw[azulFuerte!60, dashed, thick] (2,-0.55) -- (2,0.55);
    \fill[azulFuerte] (1,0) circle (2.5pt);
    \fill[azulFuerte!60] (2,0) circle (2.5pt);
    \node[azulFuerte] at (1, -0.55) {\small$1$};
    \node[azulFuerte] at (2, -0.55) {\small$2$};
    \node[azulFuerte] at (1.5,-0.9) {\small Dominio $[1,2]$};
\end{tikzpicture}
\end{center}

\[
    q(x) < 0 \quad\Longleftrightarrow\quad
    x \in \left(\frac{1-\sqrt{5}}{2},\;\frac{1+\sqrt{5}}{2}\right)
\]

\end{paso}

\bigskip

% ============================================================
% PASO 5 — INTERSECCIÓN CON EL DOMINIO
% ============================================================
\begin{paso}{5: Intersección con el Dominio}

Debemos tomar la intersección entre la solución de $q(x) < 0$ y el dominio $[1,2]$:

\[
    \underbrace{\left(\frac{1-\sqrt{5}}{2},\;\frac{1+\sqrt{5}}{2}\right)}_{\text{solución } q<0}
    \;\bigcap\;
    \underbrace{[\,1,\,2\,]}_{\text{dominio}}
    \;=\;
    \left[\,1,\;\frac{1+\sqrt{5}}{2}\right)
\]

\medskip
Nótese que:
\begin{itemize}[leftmargin=1.5cm]
    \item $x = 1$: pertenece al dominio y satisface $q(1) = 1-1-1 = -1 < 0$. \checkmark
    \item $x = \dfrac{1+\sqrt{5}}{2} \approx 1.618$: raíz de $q$, por lo tanto $q = 0$,
          no cumple la desigualdad estricta. \textbf{Se excluye.}
    \item $x = 2$: $q(2) = 4-2-1 = 1 > 0$. No cumple. \textbf{Se excluye.}
\end{itemize}

\end{paso}

\bigskip

% ============================================================
% RESULTADO FINAL
% ============================================================
\begin{tcolorbox}[colback=azulFuerte, colframe=azulFuerte,
    colupper=white, arc=10pt, boxrule=1.5pt, enhanced,
    drop shadow={black!15}]
    \centering
    \large\textbf{RESULTADO FINAL}\\[0.5em]
    \LARGE
    $\displaystyle
        x \;\in\; \left[\,1,\;\frac{1+\sqrt{5}}{2}\right)
        \;\approx\; [\,1,\;1.618\,)$
\end{tcolorbox}

\vspace{1cm}
\begin{center}
    \small\textbf{Usach Premium} — ¡Nos vemos en la ayudantía! \\
    \textit{"Por y para estudiantes."}
\end{center}

\end{document}
